% !TEX root = ../main.tex
\chapter{等式推論とリファクタリング}
\label{chap:ch09_equality_refactoring}
\BookIndex{りふぁくたりんぐ}{リファクタリング}
\BookIndex{とうしきすいろん}{等式推論}

\begin{chapterabstract}
本章では、リファクタリングを「意味を保存する等式」として扱います。
ソフトウェアでは、関数の分割、共通処理の抽出、計算式の整理、処理順序の入れ替えといった変更が日常的に行われます。
これらの変更が安全であるとは、少なくとも観測したい範囲で、変更前と変更後が同じ結果を返すということです。
Lean 4 では、この主張を等式として書き、\texttt{rfl}、\texttt{simp}、\texttt{rw}、\texttt{calc} で確認できます。

この章では、Lean で大きなプログラム全体を検証することを目的にしません。
変更で壊したくない性質を小さな式として切り出し、その式がすべての入力で成り立つことを確認する方法を学びます。
\end{chapterabstract}

\begin{learninggoals}
\item リファクタリングを「変更前後の意味保存」として説明できます。
\item Lean で、二つの関数が同じ結果を返すことを等式として書けます。
\item \texttt{simp}、\texttt{rw}、\texttt{calc} を、式の整理、書き換え、レビュー可能な変形ログとして読めます。
\item 処理順序を変えてよい場合と、変えてはいけない場合を、証明可能な等式によって区別できます。
\end{learninggoals}

\section{全入力で振る舞いを保存する}

リファクタリングは、テストを書くだけで常に安全になるわけではありません。
テストは、選んだ入力に対して結果が変わらないことを確認します。
しかし、リファクタリングで述べたい性質は、通常は次の形です。

\begin{quote}
すべての入力について、変更前の処理と変更後の処理は同じ結果を返す。
\end{quote}

これは、自然言語のレビューコメントとしても書けます。
Lean では、この主張をそのまま命題として書けます。
たとえば、\texttt{old x = new x} という等式を、任意の \texttt{x} について証明します。

図\ref{fig:ch09-equality-as-observation}は、実装の構造そのものではなく、同じ入力に対する戻り値を観測して変更前後を比較します。

\FigurePlaceholder{fig-ch09-equality-as-observation.png}{0.90}{変更前後の観測}{fig:ch09-equality-as-observation}

\section{等しいとは何か}

\subsection{見た目が同じことと、意味が同じこと}

ソフトウェアでは、二つのコードが同じであるという表現に複数の意味があります。
文字列として同じ、ASTとして同じ、実行速度が同じ、外部から見た結果が同じ、ログや副作用まで同じ、などの意味です。

本章では、最初の一歩として「戻り値の等しさ」を扱います。
つまり、同じ入力を与えたときに同じ値を返すという性質です。
これは、純粋な関数を小さくモデル化するときに扱いやすい性質です。

本章でいう意味保存とは、変更前の関数 \texttt{old} と変更後の関数 \texttt{new} について、すべての入力 \texttt{x} で \texttt{old x = new x} が成り立つことです。

ただし、戻り値の等しさは、現実のリファクタリング安全性のすべてではありません。
実システムでは、例外、ログ、DB更新、時刻、乱数、外部API呼び出し、パフォーマンスなども観測されることがあります。
本章では、まず戻り値に集中します。

\subsection{Leanにおける等式}

Lean では、等式 \texttt{a = b} は命題です。
命題であるため、証明の対象になります。

次の例では、\texttt{rawCalc} は少し冗長な計算であり、\texttt{refactoredCalc} はそれを整理したものと考えます。

\begin{listing}[htbp]
\begin{minted}{lean4}
namespace Chapter09

-- 何もしない整理は、simp で消せることが多い。
def rawCalc (x : Nat) : Nat :=
  (x + 0) * 1

def refactoredCalc (x : Nat) : Nat :=
  x

theorem rawCalc_eq_refactoredCalc (x : Nat) :
    rawCalc x = refactoredCalc x := by
  unfold rawCalc refactoredCalc
  simp
\end{minted}
\caption{\texttt{rawCalc} と \texttt{refactoredCalc}}
\label{lst:ch09_raw_refactored}
\end{listing}

\texttt{namespace Chapter09} は、この章の名前を他章と衝突させないための囲いです。
定理は、\texttt{rawCalc} と \texttt{refactoredCalc} が任意の自然数 \texttt{x} で同じ結果を返すと述べます。
\texttt{unfold} で二つの定義を展開し、\texttt{simp} で \texttt{x + 0} と \texttt{x * 1} を \texttt{x} へ整理しています。

\section{関数の結果が同じであることを証明する}

\subsection{関数そのものではなく、観測結果を比較する}

実務で確認したいものは、しばしば「実装が同じか」ではなく「外から見た結果が同じか」です。
内部で一時変数を使ったり、補助関数を抽出したりしても、同じ入力に対して同じ出力を返すなら、その観測範囲では安全な変更だと言えます。

この考え方を Lean で表すには、「すべての入力について結果が等しい」という命題を書きます。

\begin{listing}[htbp]
\begin{minted}{lean4}
-- 関数の結果が同じ、という仕様を全入力についての命題にする。
def sameOnNat (f g : Nat → Nat) : Prop :=
  ∀ x, f x = g x

theorem raw_and_refactored_same :
    sameOnNat rawCalc refactoredCalc := by
  intro x
  unfold rawCalc refactoredCalc
  simp
\end{minted}
\caption{\texttt{sameOnNat}}
\label{lst:ch09_same_on_nat}
\end{listing}

\texttt{sameOnNat} は、二つの \texttt{Nat → Nat} 関数が同じ結果を返すことを表す述語です。
\texttt{intro x} は、「任意の \texttt{x} について示す」ために、入力 \texttt{x} を一つ仮に取る操作です。
その \texttt{x} について等式を証明できれば、任意の \texttt{x} について証明したことになります。

これは、テストケースを一つ増やすこととは異なります。
\texttt{intro x} の \texttt{x} は特定の値ではなく任意の入力を表すため、証明はモデル化した範囲の全入力を覆います。

\subsection{関数外延性との関係}

数学では、二つの関数が同じであることを「すべての入力で同じ出力を返すこと」として扱います。
この考え方を関数外延性と呼びます。
\texttt{sameOnNat} は、この考え方を \texttt{∀ x, old x = new x} という点ごとの等式で表します。
この形は、実務でも「すべての入力 \texttt{x} で結果が一致する」とそのまま読めます。

\section{\texttt{calc} による式変形}

\subsection{証明をレビュー可能なログにする}

\texttt{simp} は便利ですが、何が起きたのかを読者に説明したい場合は、変形を段階的に書く方が明確です。
そのために \texttt{calc} を使います。

\texttt{calc} は、式を一行ずつ変形するための構文です。
コードレビューで「まず括弧を付け替え、次に補助変数の形に合わせる」と説明することに似ています。

ここでは、料金計算を例にします。
変更前は式をそのまま書いています。
変更後は、小計 \texttt{subtotal} と追加費用 \texttt{extra} を一時変数として抽出しています。
\texttt{let x := v} は、その式の中だけで値 \texttt{v} に名前 \texttt{x} を付ける局所定義です。

\begin{listing}[htbp]
\begin{minted}{lean4}
-- calc は、レビューしやすい段階的な式変形である。
def totalOriginal (price qty shipping handling : Nat) : Nat :=
  ((price * qty) + shipping) + handling

def totalExtracted (price qty shipping handling : Nat) : Nat :=
  let subtotal := price * qty
  let extra := shipping + handling
  subtotal + extra

theorem total_extracted_preserves
    (price qty shipping handling : Nat) :
    totalOriginal price qty shipping handling =
      totalExtracted price qty shipping handling := by
  unfold totalOriginal totalExtracted
  calc
    ((price * qty) + shipping) + handling
        = (price * qty) + (shipping + handling) := by
          rw [Nat.add_assoc]
\end{minted}
\caption{\texttt{total\_extracted\_preserves}}
\label{lst:ch09_calc_extract}
\end{listing}

\texttt{calc} の各行は、直前の式から次の式へ進む根拠を示します。
この証明では、足し算の結合律 \texttt{(a + b) + c = a + (b + c)} を使います。
これにより、\texttt{shipping} と \texttt{handling} を先にまとめても合計が変わらないことを確認しています。

\subsection{小さな法則を、リファクタリング規則として読む}

数学の定理名は、初めは難しく見えることがあります。
しかし、実務上は次のように読み替えられます。

\begin{itemize}
\item \texttt{Nat.add\_assoc}：足し算の括弧を付け替えてもかまいません。
\item \texttt{Nat.add\_comm}：足し算の左右を入れ替えてもかまいません。
\item \texttt{simp}：0を足す、1を掛ける、定義上明らかな分岐を整理します。
\end{itemize}

Lean の定理は、暗記する対象というより、リファクタリングで使える安全な書き換え規則です。

結合律や交換律は、どの処理にも成り立つわけではありません。
足し算では順序を入れ替えられますが、文字列連結、丸めを含む料金計算、ログ出力、DB更新、エラーを返す処理では、順序を変えると結果や観測が変わることがあります。

\section{書き換えとしてのリファクタリング}

\subsection{\texttt{rw} は仕様に基づく置換である}

\texttt{rw} は rewrite の略です。
すでに証明済みの等式を使って、現在のゴールを書き換えます。

次の例では、\texttt{price + fee} を \texttt{fee + price} に入れ替えます。
自然数の足し算では、この入れ替えは安全です。
Lean では、その根拠として \texttt{Nat.add\_comm} を使います。

\begin{listing}[htbp]
\begin{minted}{lean4}
-- rw は、既知の等式を使った置換として読める。
def addFeeOriginal (price fee : Nat) : Nat :=
  price + fee

def addFeeRefactored (price fee : Nat) : Nat :=
  fee + price

theorem addFee_refactor_preserves (price fee : Nat) :
    addFeeOriginal price fee = addFeeRefactored price fee := by
  unfold addFeeOriginal addFeeRefactored
  rw [Nat.add_comm]
\end{minted}
\caption{\texttt{addFee\_refactor\_preserves}}
\label{lst:ch09_rw_comm}
\end{listing}

この証明は短いですが、根拠が明確です。
レビューでいう「足し算なので順序を入れ替えられる」という根拠が、ここでは \texttt{rw [Nat.add\_comm]} として機械的に検査されます。

\subsection{補題を使って証明を短くする}

実務の証明では、毎回大きな式を直接変形すると読みにくくなります。
その場合は、小さな補題を先に作ります。

補題とは、後で使うための小さな定理です。
たとえば、「手数料Aと手数料Bの加算順序は入れ替えてよい」という補題を作っておけば、より大きな料金計算の証明でその補題を使えます。

この章では、大きな補題の設計には立ち入りません。
リファクタリングの根拠を、名前付きの小さな証明として切り出せることを確認します。

図\ref{fig:ch09-refactoring-rewrite}は、変更前の式から変更後の式までを、根拠のある小さな等式の列としてつなぎます。

\FigurePlaceholder{fig-ch09-refactoring-rewrite.png}{0.90}{書き換え規則の列}{fig:ch09-refactoring-rewrite}

\section{サンプル：処理順序を変えたときの同値性}

\subsection{順序変更は、いつでも安全ではない}

処理順序を変えるリファクタリングには注意が必要です。
特に、丸め、上限・下限、失敗、状態更新、副作用を含む場合は、順序が意味を持ちます。

ただし、一定の条件では順序を変えられます。
たとえば、二つの処理がどちらも単なる加算であれば、順序を入れ替えても最終的な合計は変わりません。
次のコードは、この性質を Lean で証明します。

\begin{listing}[htbp]
\begin{minted}{lean4}
-- 処理順序を変えてよいのは、必要な等式が証明できる場合だけである。
def addShipping (n : Nat) : Nat :=
  n + 500

def addHandling (n : Nat) : Nat :=
  n + 100

theorem add_steps_commute (n : Nat) :
    addHandling (addShipping n) = addShipping (addHandling n) := by
  unfold addHandling addShipping
  calc
    (n + 500) + 100 = n + (500 + 100) := by
      rw [Nat.add_assoc]
    _ = n + (100 + 500) := by
      rw [Nat.add_comm 500 100]
    _ = (n + 100) + 500 := by
      rw [← Nat.add_assoc]
\end{minted}
\caption{\texttt{add\_steps\_commute}}
\label{lst:ch09_order_commutes}
\end{listing}

この証明は、次の三段階で読めます。

\begin{enumerate}
\item まず、\texttt{(n + 500) + 100} の括弧を付け替えます。
\item 次に、\texttt{500 + 100} を \texttt{100 + 500} に入れ替えます。
\item 最後に、括弧を戻して \texttt{(n + 100) + 500} の形にします。
\end{enumerate}

これは、処理順序を入れ替えるためのレビュー可能な変形ログです。

\subsection{順序を変えると壊れる例}

一方、割引のような処理は、順序によって結果が変わることがあります。
次の例では、\texttt{Nat} の引き算を使います。
Lean の \texttt{Nat} では、引き算の結果は負の値にはならず、0で止まります。
この単純化されたモデルでも、順序の違いを確認できます。

\begin{listing}[htbp]
\begin{minted}{lean4}
-- 反例の雰囲気: Nat の引き算は 0 で止まるため、順序で結果が変わる。
def discount100 (n : Nat) : Nat :=
  n - 100

#eval discount100 (addShipping 50)  -- => 450
#eval addShipping (discount100 50)  -- => 500
\end{minted}
\caption{\texttt{discount100}・\texttt{addShipping}}
\label{lst:ch09_order_not_safe}
\end{listing}

最初の式は、送料を足してから割引します。
二つ目の式は、割引してから送料を足します。
\texttt{50} という入力では結果が変わるため、この二つの処理は一般には入れ替えられません。

処理が似ているという理由だけで、順序を変えてはいけません。
順序変更を安全だと主張するには、対応する等式を証明する必要があります。
証明できない場合は、反例があるか、モデルに必要な前提条件が不足している可能性があります。

\section{実務での切り出し方}

Lean に実務コードをそのまま持ち込む必要はありません。
最初は、持ち込まない方が扱いやすい場合もあります。
変更で壊したくない性質を、小さく切り出します。

たとえば、大きな料金計算の実装があるとします。
その全体には、通貨、丸め、税率、キャンペーン、ユーザー種別、外部設定などが関わる可能性があります。
最初からすべてを形式化すると、学習コストと保守コストが高くなります。

最初に対象とするものは、次のような小さな核です。

\begin{itemize}
\item 0を足す処理を消しても結果は変わりません。
\item 同じ種類の加算処理をまとめても結果は変わりません。
\item 補助関数を抽出しても、外から見た戻り値は変わりません。
\item ある二つの処理は、特定の前提のもとで順序を入れ替えられます。
\item ある二つの処理は、順序を入れ替えてはいけません。
\end{itemize}

\section{よくある誤解}

等式で証明しても、実システム全体が安全になるわけではありません。
Lean で証明したものは、Lean 上に書いたモデルと命題です。
本番コード、外部API、DB、時刻、乱数、並行処理などが関わる場合は、どこまでをモデル化したのかを明示する必要があります。

\subsection{誤解1：テストは不要になる}

Lean の証明を導入しても、テストは不要になりません。
Lean の証明は、モデル化した範囲で全称的な保証を与えます。
一方、テストは、本番コード、I/O、外部サービス、パフォーマンス、設定ファイル、環境依存の挙動を確認する場合に適しています。

本章で扱った証明は、テストの代替ではなく、テストでは網羅しにくい性質を補うために使います。

\subsection{誤解2：すべてのリファクタリングを証明すべきである}

すべてのリファクタリングを証明する必要はありません。
小さすぎる変更や、影響が限定的な変更まで証明すると、形式化のコストが価値を上回ることがあります。
証明の対象には、壊れると困る仕様の核を選びます。

\subsection{誤解3：\texttt{simp} が通れば意味を理解しなくてよい}

\texttt{simp} は便利ですが、何が単純化されたのかを人間が理解していないと、証明をレビュー資料として使いにくくなります。
初学者のうちは、\texttt{calc} を使って変形を明示する練習も役立ちます。

\section{本章のまとめ}

要点と記法
\begin{itemize}
\item 観測対象を戻り値に限る場合、意味保存は \texttt{old x = new x} と書けます。
\item 順序変更には、結合律や交換律など、その等式を正当化する条件が必要です。
\item \texttt{simp} / \texttt{rw} は、単純化と、証明済み等式による書き換えに使います。
\item \texttt{calc} は、等式や不等式の変形を段階的に並べるブロックです。
\item \texttt{rw [← h]} は、等式 \texttt{h} を逆向きに使って書き換えます。
\end{itemize}
